% 原书分章PDF第3—6页：第30.2节全部10段。
\section{认识论与方法论中的真理}

% S30.2.P01
为说明真理的这种困境，或者更确切地说，为理解它目前的处境，我们可以考虑采取历史哲学的立场。这里不探讨真理的种种表现形式，而将范围限于一般的思想史，尤其是哲学——特别是形而上学与认识论——以及现代科学。这样一来，我们会看到贯穿各个时代的一种根本张力：一方是谦逊而无私的“真理的仆人”（ancillae veritatis），他们追求真知识这一圣杯；另一方则认为，这一切不过是乌托邦，要么因为真理概念本身定义不清或构思不当，要么仅仅因为人类自身的处境使我们无法实现它。

% S30.2.P02
例如，这种张力很容易追溯到古希腊，起点是柏拉图（Plato）的 episteme 概念，即作为“得到证成的真信念”的真实、确定的知识，以及亚里士多德（Aristotle）的公理演绎式知识理想和建立在保真推理之上的三段论逻辑。episteme 被认为与 doxa，即意见、不确定性和不规则性相对立。这些理想几乎立即遭到智者修辞学运动——包括前述普罗泰戈拉——以及第欧根尼（Diogenes）的犬儒学派的挑战与反抗；而最为显著、也最为深刻的挑战来自怀疑论者。尤其是埃利斯的皮浪（Pyrrho of Elis），为批判关于真理的主张铺平了道路；塞克斯都·恩披里柯（Sextus Empiricus）在公元2世纪或3世纪的《皮浪主义原理》（Principles of Pyrrhonism）中分析并传递了他的思考。怀疑论以各种面貌和形式，像波浪一样贯穿历史：从16世纪米歇尔·德·蒙田（Michel de Montaigne）的怀疑论人文主义，经17世纪勒内·笛卡尔（Rene Descartes）的方法论怀疑、18世纪大卫·休谟（David Hume），一直延续到20世纪卡尔·波普尔（Karl Popper）及当代哲学。上述张力持续至今，因此，这一理想即使在大数据与数据科学时代，也依然具有意义。不过，随着哲学的演变和科学的进步，对真理的关注至少沿两条不同路线继续发展。这里将作简要阐述：第一条在认识论／科学哲学之中，第二条在研究方法论之中。

% S30.2.P03
科学哲学（philosophy of science）经历了许多发展，它们影响或塑造了真理目前的地位；本文只突出其中几个重要方面。首先，一个根本问题是：真理究竟是实在、某种状态、过程、事实、思想，还是命题／陈述／理论／模型的属性？这一问题催生了多种真理的概念化方式或定义，并在认识论中得到发展和审视。它们包括“严格”的数学版本，例如真理符合论（correspondence theory of truth）及阿尔弗雷德·塔斯基（Alfred Tarski）与之相关的真理语义论（semantic theory of truth）；“温和”的版本，例如真理融贯论（coherence theory of truth）；以及或多或少带有“相对主义”色彩的版本，例如真理共识论（consensus theory of truth）、约定论方法或实用主义方法。

% S30.2.P04
确实，人们常常主张，真正具有数学精确性的真理概念，只能在逻辑与模型论中成立，而不能在经验科学中成立。数理逻辑具有真值函数性质（truth-functional）；逻辑蕴涵、推理与有效性等关键概念，都以这种方式理解，而逻辑可靠性（soundness）与完全性（completeness）准则，则指导并保证推理过程的正确性。如果一个公式可以被推导出来，它就必须为真，这就是可靠性；如果它为真，它就应当能够被推导出来，这就是完全性。这种把句法可推导性与语义真视为对偶概念的严格符合论，在经验科学的认识论中几乎没有引起共鸣，尽管可以认为，方法论中的某些著名应用能够追溯到它。

% S30.2.P05
例如，在流行病学中，敏感度（sensitivity）、特异度（specificity），以及由它们衍生的概念——包括阳性与阴性检验的似然比——以类似方式定义检验后概率。在信息检索与模式识别中，查准率（precision）和查全率（recall）也有相应定义。前者表达的意思是：如果一份文献被检索出来（被推导出来），它就应当是相关的（真的）。反过来，查全率表达的是：如果一篇文本是相关的，我们就希望它被检索出来。显然，这些都是用于研究方法论的概率概念，而不是认识论本身的支柱。后面将更全面地讨论这一主题。

% S30.2.P06
其次，自20世纪初以来，在盎格鲁—撒克逊的分析哲学传统中，莫里茨·石里克（Moritz Schlick）、Rudolph Carnap 等逻辑实证主义者（logical positivists），通过可证实性标准与归纳确证，将真理概念与意义联系起来；后来，它又与证伪主义相联系，主要体现在波普尔的批判理性主义（critical rationalism）之中。然而，所有这些卓越学者所关注的都是证成语境（context of justification），而不是发现语境（context of discovery）。因此，他们主要旨在对经验科学作理性重构，并运用概率概念实现这一目标。他们并不致力于发展一种建设性的研究方法论。而且，真理实际上是在更广泛的认识论问题背景中被研究的，例如客观性、理性、科学理论的结构、价值中立性、科学的统一性等。

% S30.2.P07
这些研究所依据的真理概念，往往建立在融贯性或实用主义之上。查尔斯·S.皮尔士（Charles S. Peirce）、威廉·詹姆斯（William James）、约翰·杜威（John Dewey）和 W.V.O. Quine 等哲学家尤其倡导后一种观念；他们都是实用主义的先驱，而实用主义至今仍是科学哲学中的主导思潮之一。他们试图使这个老式概念更适合当代认识论问题，也更能在其中成立。逻辑实证主义关于科学理论结构的研究，同时导向了句法方法（Carnap）与语义方法；后者把理论描述为模型的类（Suppes、Van Fraassen、Stegmüller）。这种区分仅仅松散地建立在上述“对偶”的逻辑符合论之上。因此，语义方法最适合归入融贯的真理观；它也似乎成为乔治·博克斯及物理学家约翰·冯·诺依曼（John von Neumann）等人当前立场的先声。这正是下一节将要讨论的问题。

% S30.2.P08
科学哲学中挑战真知识理想的第三项发展，是具有影响力的历史学与社会学研究的兴起，例如托马斯·库恩（Thomas Kuhn）、保罗·费耶阿本德（Paul Feyerabend）和理查德·罗蒂（Richard Rorty）的研究，法兰克福学派（Frankfurter Schule）的马克思主义研究，Robert K. Merton、Bruno Latour 发表的外在论研究，David Bloor 和 Barry Barnes 的“强纲领”（Strong Program），以及较近时期 Samuel Pickering 的社会建构主义。例如，库恩在《科学革命的结构》（The Structure of Scientific Revolutions，1962）中，提出了著名的范式理论：科学发展表现为不连续，而非连续进步；不同范式所使用的概念和方法不可通约；观察负载着理论，等等。保罗·费耶阿本德在《反对方法》（Against Method，1975）中，以其无政府主义观念和“怎么都行”（anything goes）的辩护震动了研究界。理查德·罗蒂在《哲学与自然之镜》（Philosophy and the Mirror of Nature，1979）中宣告了表征的终结，即否定科学理论以某种方式模仿自然这一观念。Bruno Latour 是网络—行动者理论（network-actor theory）的奠基者之一；他详细研究了实验科学家实际上如何在实验室中开展科学工作，包括做实验、讨论、协商、与利益攸关方及资助机构互动，以及处理惯例与冲突。他的《科学在行动：怎样在社会中跟随科学家与工程师》（Science in Action: How to Follow Scientists and Engineers through Society，1987）是这一领域的现代经典，并影响了许多科学技术研究。

% S30.2.P09
尽管这些学者都把更广泛意义上的整个科学事业作为研究对象，包括客观性、理性，以及科学被宣称或受到争议的价值中立地位，他们的研究仍然导向了各类相对主义，也影响了真理观念。无论源于外在论的历史研究、当代怀疑论、文化相对主义，还是社会建构主义，真理概念都受到了侵蚀；这正是今日状况的特征。虽然真理是认识论的关键概念，但颇具悖论意味的是，人们往往把它当作一个初始的或不言自明的术语，隐含地、默默地使用它，并在对使用者有利时诉诸它；而认识论真正的研究重点，却是前文已提到的那些不那么根本的概念。因此，与真理概念相关的一系列问题，常常以间接方式显现，并持续出现在当代科学哲学的议题中。它们不仅包含或涉及规范哲学中的传统主题——伦理、价值、诚信、道德责任、知情同意、行为准则等——而且真理概念尤其会在经典的认识论问题中出现、发挥作用或被预设，例如科学实在论之争（scientific realism debate）、（贝叶斯）确证理论、科学理论的结构，以及最后但同样重要的“解释、因果性与规律”这一三联结构；这些概念往往被一并研究，或被置于彼此联系之中。

% S30.2.P10
关于真理的思考沿之发展的第二条路线，涉及研究实践与方法论的演变。粗略地说，真理概念已经被操作化、被重塑，或者更确切地说，由效度（validity）、信度（reliability）、可靠性（soundness）、正确性或准确性等概念加以补充。例如，研究方法论已经运用了大量效度概念，包括内部与外部效度、构念效度、统计结论效度、同时效度；这个清单很容易继续延伸。信度也同样如此，它具有许多内涵，构成了真理及其表征与“测量”的前提或基础。此外，人们还发展了许多或多或少依赖具体领域的质量度量或标准，包括前述敏感度、特异度、查全率、查准率等概念。尽管这两方面的发展都很重要，它们却未能阻止前述局面的出现；而且，人们过于频繁地以不恰当的权威姿态诉诸某种真理观念，其结果是在利用真理的脆弱性，而不是保护真理。具有讽刺意味的是，一方面，这个概念似乎已经过时，被绕开、被忽视；另一方面，我们似乎又经常诉诸它，需要它，并预设它。这里将说明，本小节如何引出若干重要侧面，从而塑造或至少照亮真理的困境。
